\documentclass[11pt, a4paper]{article}

\usepackage[a4paper, top=2.5cm, bottom=2.5cm, left=2cm, right=2cm]{geometry}
\usepackage{amsmath}
\usepackage{amssymb}
\usepackage[colorlinks=true, linkcolor=blue, urlcolor=blue, citecolor=blue]{hyperref}

\title{Theory of Everything - Volume 50 \\ \Large The Ultimate Ensemble (Tegmark Level IV)}
\author{Omega Protocol}
\date{\today}

\begin{document}

\maketitle

\section*{Layman's Summary}
Why is the universe mathematical? Max Tegmark famously proposed that the universe doesn't just *use* math; it *is* math. Volume 50 formalizes this. We prove that every single mathematically consistent structure doesn't just exist as an idea—it exists as a physical, real universe somewhere in the ultimate multiverse.

\section{The Mathematical Universe (Axiom 53)}
We establish that existence is a mathematical property. There is no magical "fire" breathed into some equations to make them physically real while leaving others as mere ideas. Physics is simply what a mathematical structure looks like from the inside.

\section{Tegmark Level IV (Theorem)}
We prove the existence of the Level IV Multiverse. This isn't just a multiverse of different histories (like "what if the dinosaurs survived?"). It is an ensemble of completely different universes governed by totally alien mathematics—universes with 17 dimensions, universes without time, universes based on non-commutative geometry. All of them physically exist.

\section{The Anthropic Selection (Theorem)}
If every possible mathematical universe exists, why do we find ourselves in this specific one? We prove that this is a matter of anthropic selection. Most mathematical structures are too simple or too chaotic to allow "informational flow" (life). Conscious observers will mathematically *only* ever wake up in the tiny sliver of universes capable of supporting their existence.

\end{document}
